\begin{table}[htbp]
\centering
\small
\caption{Problema A (google/gemini-2.5-flash-lite): métrica estimada em zero-shot, few-shot e pelo rótulo real.}
\label{tab:T1_A_flashlite}
\begin{tabular}{lrrrrlrrl}
\toprule
 &  &  & \multicolumn{3}{c}{Perceptron} & \multicolumn{3}{c}{NNLS} \\ \cmidrule(lr){4-6} \cmidrule(lr){7-9}
Regime & Acur.\ LLM vs real & Partição LLM (n0/n1) & Fid.\ (vs LLM) & Acur.\ vs real & $\hat{W}$ unit. & Fid.\ (vs LLM) & Acur.\ vs real & $\hat{W}$ unit. \\
\midrule
zero-shot† & 25,1 (21,3–29,3) & 46/104 & 81,3 (78,7–83,3) & 17,6 (15,3–21,3) & [0,27; 0,96] & 82,7 (80,7–84,0) & 17,1 (14,7–19,3) & [0,26; 0,96] \\
4-shot† & 88,7 (57,5–98,6) & 78/68 & 93,8 (63,0–99,3) & 92,6 (83,3–98,0) & [0,43; 0,86] & 93,3 (59,6–100,0) & 92,1 (73,3–98,0) & [0,45; 0,73] \\
10-shot† & 94,0 (87,1–97,9) & 71/69 & 97,5 (87,1–100,0) & 94,7 (88,7–98,0) & [0,61; 0,74] & 97,5 (88,6–100,0) & 95,1 (92,7–98,0) & [0,67; 0,50] \\
20-shot† & 94,1 (88,5–97,7) & 66/64 & 97,9 (95,4–100,0) & 95,4 (92,7–98,0) & [0,67; 0,70] & 98,1 (95,4–100,0) & 95,6 (92,7–98,0) & [0,69; 0,50] \\
40-shot† & 95,1 (90,9–98,2) & 55/55 & 99,0 (95,5–100,0) & 95,3 (92,7–98,0) & [0,70; 0,71] & 99,0 (95,5–100,0) & 95,3 (92,7–98,0) & [0,81; 0,39] \\
\midrule
W ajustado ao rótulo real (teto)† & — & — & — & 96,2 (95,3–98,0) & [0,37; 0,89] & — & 95,3 (93,3–98,0) & [0,51; 0,86] \\
\bottomrule
\multicolumn{9}{p{0.95\linewidth}}{\footnotesize Médias das sementes e repetições, mín–máx entre parênteses, em \%. Fid.\ (vs LLM) = fidelidade: fração dos pontos em que a métrica reproduz o rótulo do LLM. Acur.\ vs real = acurácia (vs rótulo real) da métrica aplicada aos 150 pontos; Acur.\ LLM vs real = acurácia (vs rótulo real) do próprio LLM, mostrada como contexto. Few-shot: exemplos rotulados pelo rótulo real (k/2 por classe, sorteio aleatório), LLM classifica os $150-k$ restantes; fidelidade medida nesses pontos. Teto: W ajustado diretamente ao rótulo real; não tem fidelidade (não é estimado a partir do LLM), só acurácia (vs rótulo real), que não chega a 100\% pela sobreposição das gaussianas do A. Em zero-shot (e com nomes de classe neutros) a correspondência entre nome e classe é arbitrária: a acurácia (vs rótulo real) do LLM pode sair abaixo de 50\%; o valor invariante à permutação é max($a$, 1 $-$ $a$). † = em pelo menos uma estimativa o Perceptron devolveu $\hat{W}$ em melhor esforço ($\gamma = 0$, sem margem viável). ``—'' = colapso em classe única (ou quase-colapso: classe minoritária com menos de 5 pontos); na linha colapsada também a acurácia do LLM e a partição ficam como ``—''.} \\
\end{tabular}
\end{table}
